\chapter{Chapter06 — Almost K Wise Algorithm}
%%% generated by the blueprint pipeline from FABL.Chapter06.Constructions.AlmostKWiseAlgorithm
\section{Overview}

Declarations of the Lean module \texttt{FABL.Chapter06.Constructions.AlmostKWiseAlgorithm}.

\section{Declarations}

\begin{definition}[Almost-wise independent input parameters]
\lean{FABL.AlmostKWiseInput}
\statementsource{Analysis-of-Boolean-Functions-by-Ryan-ODonnell}{
  pdf-ffbe9c7b3306-p155-b007,
  pdf-ffbe9c7b3306-p170-b003
}
\label{FABL.AlmostKWiseInput}
\leanok
An almost $k$-wise input parameter record consists of natural numbers $k,n,a,b$ such
that $1\le k\le n$, $a>0$, and $2a\le b$, where $a/b$ is the specified rational bias.
\end{definition}

\begin{theorem}[Positivity of the output dimension]
\lean{FABL.AlmostKWiseInput.n_pos}
\label{FABL.AlmostKWiseInput.n_pos}
\leanok
\uses{FABL.AlmostKWiseInput}
\difficulty{2}
Let $k,n,N,D$ be nonnegative integers satisfying $1\le k\le n$, $0<N$, and $2N\le D$.
Then $n>0$.
\end{theorem}

\begin{definition}[Vandermonde input parameters]
\lean{FABL.AlmostKWiseInput.vandermondeInput}
\statementsource{Analysis-of-Boolean-Functions-by-Ryan-ODonnell}{
  pdf-ffbe9c7b3306-p154-b007,
  pdf-ffbe9c7b3306-p170-b003
}
\label{FABL.AlmostKWiseInput.vandermondeInput}
\leanok
\uses{FABL.AlmostKWiseInput, FABL.ExecutableVandermondeInput}
Given an almost $k$-wise input with independence order $k$ and output dimension $n$,
form the corresponding Vandermonde input with the same $k$ and $n$.
\end{definition}

\begin{definition}[Row count]
\lean{FABL.AlmostKWiseInput.rowCount}
\statementsource{Analysis-of-Boolean-Functions-by-Ryan-ODonnell}{
  pdf-ffbe9c7b3306-p154-b006,
  pdf-ffbe9c7b3306-p170-b003
}
\label{FABL.AlmostKWiseInput.rowCount}
\leanok
\uses{FABL.AlmostKWiseInput, FABL.AlmostKWiseInput.vandermondeInput, FABL.ExecutableVandermondeInput.fieldDegree}
For an almost $k$-wise input, let $d$ be the field degree associated with its
Vandermonde input. Define the row count to be $(k-1)d+1$.
\end{definition}

\begin{theorem}[Positivity of the row count]
\lean{FABL.AlmostKWiseInput.rowCount_pos}
\label{FABL.AlmostKWiseInput.rowCount_pos}
\leanok
\uses{FABL.AlmostKWiseInput, FABL.AlmostKWiseInput.rowCount}
\difficulty{1}
Let $k,n,N,D$ be nonnegative integers satisfying $1\le k\le n$, $0<N$, and $2N\le D$.
Define $d=\max(1,\lceil\log_2 n\rceil)$ and $r=(k-1)d+1$. Then $r>0$.
\end{theorem}

\begin{definition}[Small-bias input parameters]
\lean{FABL.AlmostKWiseInput.smallBiasInput}
\statementsource{Analysis-of-Boolean-Functions-by-Ryan-ODonnell}{
  pdf-ffbe9c7b3306-p153-b002,
  pdf-ffbe9c7b3306-p170-b003
}
\label{FABL.AlmostKWiseInput.smallBiasInput}
\leanok
\uses{FABL.AlmostKWiseInput, FABL.AlmostKWiseInput.rowCount, FABL.AlmostKWiseInput.rowCount_pos, FABL.SmallBiasInput}
Given an almost $k$-wise input, form a small-bias input whose dimension is its row count
and whose bias numerator and denominator are those of the original input.
\end{definition}

\begin{definition}[Encoded bias]
\lean{FABL.AlmostKWiseInput.epsilon}
\statementsource{Analysis-of-Boolean-Functions-by-Ryan-ODonnell}{
  pdf-ffbe9c7b3306-p155-b007,
  pdf-ffbe9c7b3306-p170-b003
}
\label{FABL.AlmostKWiseInput.epsilon}
\leanok
\uses{FABL.AlmostKWiseInput, FABL.AlmostKWiseInput.smallBiasInput, FABL.SmallBiasInput.epsilon}
For an almost $k$-wise input with bias numerator $a$ and denominator $b$, define its
real bias parameter to be $a/b$.
\end{definition}

\begin{theorem}[Positivity of the bias]
\lean{FABL.AlmostKWiseInput.epsilon_pos}
\label{FABL.AlmostKWiseInput.epsilon_pos}
\leanok
\uses{FABL.AlmostKWiseInput, FABL.AlmostKWiseInput.epsilon, FABL.AlmostKWiseInput.smallBiasInput, FABL.SmallBiasInput.epsilon_pos}
\difficulty{1}
Let $k,n,a,d$ be natural numbers with $1\leq k\leq n$, $a>0$, and $2a\leq d$. For the
bias $\varepsilon=a/d$ associated with this input, one has $\varepsilon>0$.
\end{theorem}

\begin{theorem}[Upper bound on the bias]
\lean{FABL.AlmostKWiseInput.epsilon_le_half}
\label{FABL.AlmostKWiseInput.epsilon_le_half}
\leanok
\uses{FABL.AlmostKWiseInput, FABL.AlmostKWiseInput.epsilon, FABL.AlmostKWiseInput.smallBiasInput, FABL.SmallBiasInput.epsilon_le_half}
\difficulty{1}
Let $k,n,a,d$ be natural numbers with $1\leq k\leq n$, $a>0$, and $2a\leq d$. For the
bias $\varepsilon=a/d$ associated with this input, one has $\varepsilon\leq\frac12$.
\end{theorem}

\begin{definition}[Size scale]
\lean{FABL.AlmostKWiseInput.sizeScale}
\statementsource{Analysis-of-Boolean-Functions-by-Ryan-ODonnell}{
  pdf-ffbe9c7b3306-p155-b007,
  pdf-ffbe9c7b3306-p170-b003
}
\label{FABL.AlmostKWiseInput.sizeScale}
\leanok
\uses{FABL.AlmostKWiseInput, FABL.AlmostKWiseInput.vandermondeInput, FABL.ExecutableVandermondeInput.fieldDegree}
For an almost $k$-wise input, let $d$ be the field degree associated with its
Vandermonde input. Define its size scale to be $k(d+1)$.
\end{definition}

\begin{theorem}[Row count bounded by the size scale]
\lean{FABL.AlmostKWiseInput.rowCount_le_sizeScale}
\label{FABL.AlmostKWiseInput.rowCount_le_sizeScale}
\leanok
\uses{FABL.AlmostKWiseInput, FABL.AlmostKWiseInput.rowCount, FABL.AlmostKWiseInput.sizeScale, FABL.AlmostKWiseInput.vandermondeInput, FABL.ExecutableVandermondeInput.fieldDegree}
\difficulty{4}
Let $k,n,N,D$ be nonnegative integers satisfying $1\le k\le n$, $0<N$, and $2N\le D$.
Define $d=\max(1,\lceil\log_2 n\rceil)$, the row count $r=(k-1)d+1$, and the size scale
$s=k(d+1)$. Then $r\le s$.
\end{theorem}

\begin{definition}[Reciprocal scale]
\lean{FABL.AlmostKWiseInput.reciprocalScale}
\statementsource{Analysis-of-Boolean-Functions-by-Ryan-ODonnell}{pdf-ffbe9c7b3306-p155-b007}
\label{FABL.AlmostKWiseInput.reciprocalScale}
\leanok
\uses{FABL.AlmostKWiseInput}
For an almost $k$-wise input with bias numerator $a$ and denominator $b$, define its
reciprocal scale to be the ceiling of $b/a$.
\end{definition}

\begin{definition}[Output scale]
\lean{FABL.AlmostKWiseInput.outputScale}
\statementsource{Analysis-of-Boolean-Functions-by-Ryan-ODonnell}{
  pdf-ffbe9c7b3306-p155-b007,
  pdf-ffbe9c7b3306-p170-b003
}
\label{FABL.AlmostKWiseInput.outputScale}
\leanok
\uses{FABL.AlmostKWiseInput, FABL.AlmostKWiseInput.reciprocalScale, FABL.AlmostKWiseInput.sizeScale}
For an almost $k$-wise input, define its output scale as the product of its size scale
and reciprocal scale.
\end{definition}

\begin{definition}[Scale]
\lean{FABL.AlmostKWiseInput.scale}
\statementsource{Analysis-of-Boolean-Functions-by-Ryan-ODonnell}{
  pdf-ffbe9c7b3306-p155-b007,
  pdf-ffbe9c7b3306-p170-b003
}
\label{FABL.AlmostKWiseInput.scale}
\leanok
\uses{FABL.AlmostKWiseInput}
For an almost $k$-wise input with output dimension $n$, bias numerator $a$, and
denominator $b$, define its scale to be the ceiling of $nb/a$.
\end{definition}

\begin{theorem}[Dimension bounded by the scale]
\lean{FABL.AlmostKWiseInput.n_le_scale}
\label{FABL.AlmostKWiseInput.n_le_scale}
\leanok
\uses{FABL.AlmostKWiseInput, FABL.AlmostKWiseInput.n_pos, FABL.AlmostKWiseInput.scale, FABL.SmallBiasInput, FABL.SmallBiasInput.n_le_scale, FABL.SmallBiasInput.scale}
\difficulty{3}
Let $k,n,a,d$ be natural numbers with $1\leq k\leq n$, $a>0$, and $2a\leq d$. Define the
scale by $\left\lceil nd/a\right\rceil$, the least natural number $m$ such that $ma\geq
nd$. Then $n$ is at most this scale.
\end{theorem}

\begin{theorem}[Small bias scale bounded by output scale]
\lean{FABL.AlmostKWiseInput.smallBiasScale_le_outputScale}
\label{FABL.AlmostKWiseInput.smallBiasScale_le_outputScale}
\leanok
\uses{FABL.AlmostKWiseInput, FABL.AlmostKWiseInput.outputScale, FABL.AlmostKWiseInput.reciprocalScale, FABL.AlmostKWiseInput.rowCount, FABL.AlmostKWiseInput.rowCount_le_sizeScale, FABL.AlmostKWiseInput.sizeScale, FABL.AlmostKWiseInput.smallBiasInput, FABL.SmallBiasInput.scale}
\difficulty{4}
Let $k,n,N,D$ be nonnegative integers satisfying $1\le k\le n$, $0<N$, and $2N\le D$.
Define $d=\max(1,\lceil\log_2 n\rceil)$, the size scale $S=k(d+1)$, and the reciprocal
scale $R=\lceil D/N\rceil$. Then the scale $\lceil nD/N\rceil$ is at most the output
scale $SR$.
\end{theorem}

\begin{theorem}[Small bias scale bound]
\lean{FABL.AlmostKWiseInput.smallBiasScale_le_four_scale_sq}
\label{FABL.AlmostKWiseInput.smallBiasScale_le_four_scale_sq}
\leanok
\uses{FABL.AlmostKWiseInput, FABL.AlmostKWiseInput.n_le_scale, FABL.AlmostKWiseInput.rowCount, FABL.AlmostKWiseInput.rowCount_le_sizeScale, FABL.AlmostKWiseInput.scale, FABL.AlmostKWiseInput.smallBiasInput, FABL.AlmostKWiseInput.vandermondeInput, FABL.ExecutableVandermondeInput.fieldDegree, FABL.ExecutableVandermondeInput.fieldDegree_succ_le, FABL.SmallBiasInput.scale}
\difficulty{5}
Let $k,n,N,D$ be nonnegative integers satisfying $1\le k\le n$, $0<N$, and $2N\le D$.
Let $s=\lceil nD/N\rceil$ be the scale computed from $n,N,D$. Then $s\le 4(s+1)^2$.
\end{theorem}

\begin{theorem}[Row count bound by the squared scale]
\lean{FABL.AlmostKWiseInput.rowCount_le_four_scale_sq}
\label{FABL.AlmostKWiseInput.rowCount_le_four_scale_sq}
\leanok
\uses{FABL.AlmostKWiseInput, FABL.AlmostKWiseInput.n_le_scale, FABL.AlmostKWiseInput.rowCount, FABL.AlmostKWiseInput.rowCount_le_sizeScale, FABL.AlmostKWiseInput.scale, FABL.AlmostKWiseInput.sizeScale, FABL.AlmostKWiseInput.vandermondeInput, FABL.ExecutableVandermondeInput.fieldDegree, FABL.ExecutableVandermondeInput.fieldDegree_succ_le}
\difficulty{4}
Let $k,n,N,D$ be nonnegative integers satisfying $1\le k\le n$, $0<N$, and $2N\le D$.
Define $d=\max(1,\lceil\log_2 n\rceil)$, $s=\lceil nD/N\rceil$, and $r=(k-1)d+1$. Then
$r\le 4(s+1)^2$.
\end{theorem}

\begin{definition}[Almost-wise construction matrix]
\lean{FABL.almostKWiseMatrix}
\statementsource{Analysis-of-Boolean-Functions-by-Ryan-ODonnell}{
  pdf-ffbe9c7b3306-p154-b006,
  pdf-ffbe9c7b3306-p154-b007,
  pdf-ffbe9c7b3306-p170-b003
}
\label{FABL.almostKWiseMatrix}
\leanok
\uses{FABL.AlmostKWiseInput, FABL.AlmostKWiseInput.rowCount, FABL.AlmostKWiseInput.vandermondeInput, FABL.executableVandermondeMatrix, FABL.𝔽₂}
Let $k,n$ be positive integers with $k\le n$, and let $\mathbb F_2$ be the field with
two elements. Set $d=\max(1,\lceil\log_2 n\rceil)$ and $r=(k-1)d+1$. The almost-wise
construction matrix is the executable Vandermonde matrix over $\mathbb F_2$ with $r$
rows and $n$ columns.
\end{definition}

\begin{definition}[Almost-wise linear transform]
\lean{FABL.almostKWiseTransform}
\statementsource{Analysis-of-Boolean-Functions-by-Ryan-ODonnell}{
  pdf-ffbe9c7b3306-p155-b004,
  pdf-ffbe9c7b3306-p170-b003
}
\label{FABL.almostKWiseTransform}
\leanok
\uses{FABL.AlmostKWiseInput, FABL.AlmostKWiseInput.rowCount, FABL.F₂Cube, FABL.almostKWiseMatrix}
Let $k,n$ be positive integers with $k\le n$, and let $\mathbb F_2$ be the field with
two elements. Set $d=\max(1,\lceil\log_2 n\rceil)$ and $r=(k-1)d+1$. The almost-wise
transform sends each vector $y\in\mathbb F_2^r$ to the row-vector product of $y$ with
the executable Vandermonde matrix over $\mathbb F_2$ having $r$ rows and $n$ columns.
\end{definition}

\begin{definition}[Almost-wise output list]
\lean{FABL.almostKWiseOutputList}
\statementsource{Analysis-of-Boolean-Functions-by-Ryan-ODonnell}{
  pdf-ffbe9c7b3306-p155-b007,
  pdf-ffbe9c7b3306-p170-b003
}
\label{FABL.almostKWiseOutputList}
\leanok
\uses{FABL.AlmostKWiseInput, FABL.AlmostKWiseInput.smallBiasInput, FABL.F₂Cube, FABL.almostKWiseTransform, FABL.deterministicSmallBiasAlgorithm}
Let $k,n$ be positive integers with $k\le n$, and let $a,b$ be positive integers with
$2a\le b$. The almost-wise output list is obtained by taking the outputs of the
deterministic small-bias algorithm for these parameters and applying the almost-wise
transform to each output, in their original order.
\end{definition}

\begin{definition}[Almost-wise output multiset]
\lean{FABL.almostKWiseOutputMultiset}
\statementsource{Analysis-of-Boolean-Functions-by-Ryan-ODonnell}{
  pdf-ffbe9c7b3306-p155-b007,
  pdf-ffbe9c7b3306-p170-b003
}
\label{FABL.almostKWiseOutputMultiset}
\leanok
\uses{FABL.AlmostKWiseInput, FABL.F₂Cube, FABL.almostKWiseOutputList}
Let $k,n$ be positive integers with $k\le n$, and let $a,b$ be positive integers with
$2a\le b$. The almost-wise output multiset consists of the elements of the almost-wise
output list, with each element counted according to its multiplicity in that list.
\end{definition}

\begin{definition}[Almost-wise output density]
\lean{FABL.almostKWiseDensity}
\statementsource{Analysis-of-Boolean-Functions-by-Ryan-ODonnell}{
  pdf-ffbe9c7b3306-p155-b007,
  pdf-ffbe9c7b3306-p170-b003
}
\label{FABL.almostKWiseDensity}
\leanok
\uses{FABL.AlmostKWiseInput, FABL.AlmostKWiseInput.smallBiasInput, FABL.ProbabilityDensity, FABL.ProbabilityDensity.pushforward, FABL.almostKWiseTransform, FABL.deterministicSmallBiasDensity}
Let $k,n$ be positive integers with $k\le n$, and let $a,b$ be positive integers with
$2a\le b$. The almost-wise output density is the pushforward, under the almost-wise
transform, of the deterministic small-bias density associated with these parameters.
\end{definition}

\begin{theorem}[Almost k wise small bias output list length]
\lean{FABL.almostKWiseSmallBiasOutputList_length}
\statementsource{Analysis-of-Boolean-Functions-by-Ryan-ODonnell}{
  pdf-ffbe9c7b3306-p153-b002,
  pdf-ffbe9c7b3306-p170-b003
}
\label{FABL.almostKWiseSmallBiasOutputList_length}
\leanok
\uses{FABL.AlmostKWiseInput, FABL.AlmostKWiseInput.smallBiasInput, FABL.SmallBiasInput.fieldDegree, FABL.deterministicSmallBiasAlgorithm, FABL.deterministicSmallBiasMultiset, FABL.deterministicSmallBiasMultiset_card}
\difficulty{3}
% dataset-pipeline: stub -- prose pending build_informal_statements.py
Statement pending regeneration.
\end{theorem}

\begin{theorem}[Length of the almost-wise output list]
\lean{FABL.almostKWiseOutputList_length}
\label{FABL.almostKWiseOutputList_length}
\leanok
\uses{FABL.AlmostKWiseInput, FABL.AlmostKWiseInput.smallBiasInput, FABL.SmallBiasInput.fieldDegree, FABL.almostKWiseOutputList, FABL.almostKWiseSmallBiasOutputList_length}
\difficulty{1}
Let an almost-wise input have independence order $k$ and output dimension $n$, with
$1\le k\le n$. Let its bias be the positive rational number whose numerator is $a$ and
denominator is $b$, where $2a\le b$. The output list has length $2^{2d}$, where $d$ is
the field degree used by the associated small-bias input.
\end{theorem}

\begin{theorem}[Output multiset cardinality preservation]
\lean{FABL.almostKWiseOutputMultiset_card_eq_source}
\label{FABL.almostKWiseOutputMultiset_card_eq_source}
\leanok
\uses{FABL.AlmostKWiseInput, FABL.AlmostKWiseInput.smallBiasInput, FABL.almostKWiseOutputList, FABL.almostKWiseOutputMultiset, FABL.deterministicSmallBiasMultiset}
\difficulty{1}
Let an almost-wise input have independence order $k$ and output dimension $n$, with
$1\le k\le n$. Let its bias be the positive rational number whose numerator is $a$ and
denominator is $b$, where $2a\le b$. The cardinality of its output multiset equals the
cardinality of the source multiset used to form it.
\end{theorem}

\begin{definition}[Almost-wise random-bit count]
\lean{FABL.almostKWiseRandomBits}
\statementsource{Analysis-of-Boolean-Functions-by-Ryan-ODonnell}{
  pdf-ffbe9c7b3306-p155-b006,
  pdf-ffbe9c7b3306-p170-b003
}
\label{FABL.almostKWiseRandomBits}
\leanok
\uses{FABL.AlmostKWiseInput, FABL.AlmostKWiseInput.smallBiasInput, FABL.SmallBiasInput.fieldDegree}
Let $k,n,a,b$ be natural numbers satisfying $1\le k\le n$, $a>0$, and $2a\le b$. The
associated random-bit count is twice the field degree of the small-bias input determined
by these parameters.
\end{definition}

\begin{theorem}[Output multiset cardinality]
\lean{FABL.almostKWiseOutputMultiset_card_eq_two_pow_randomBits}
\label{FABL.almostKWiseOutputMultiset_card_eq_two_pow_randomBits}
\leanok
\uses{FABL.AlmostKWiseInput, FABL.almostKWiseOutputList_length, FABL.almostKWiseOutputMultiset, FABL.almostKWiseRandomBits}
\difficulty{1}
Let an almost-wise input have independence order $k$ and output dimension $n$, with
$1\le k\le n$. Let its bias be the positive rational number whose numerator is $a$ and
denominator is $b$, where $2a\le b$. If $r$ is the number of random bits specified by
this input, then its output multiset has cardinality $2^r$.
\end{theorem}

\begin{theorem}[Fourier regularity of the constructed density]
\lean{FABL.almostKWiseDensity_isApproximatelyKWiseIndependent}
\proofsource{Analysis-of-Boolean-Functions-by-Ryan-ODonnell}{
  pdf-ffbe9c7b3306-p155-b004,
  pdf-ffbe9c7b3306-p155-b005
}
\label{FABL.almostKWiseDensity_isApproximatelyKWiseIndependent}
\leanok
\uses{FABL.AlmostKWiseInput, FABL.AlmostKWiseInput.epsilon, FABL.AlmostKWiseInput.smallBiasInput, FABL.AlmostKWiseInput.vandermondeInput, FABL.IsLowDegreeFourierRegular, FABL.ProbabilityDensity.IsBiased, FABL.almostKWiseDensity, FABL.almostKWiseMatrix, FABL.binaryFunctionOnSignCube, FABL.deterministicSmallBiasAlgorithm_spec, FABL.deterministicSmallBiasDensity, FABL.executableVandermondeMatrix_hasNonzeroColumnSumsUpTo, FABL.matrixPushforwardDensity, FABL.matrixPushforwardDensity_isApproximatelyKWiseIndependent}
\difficulty{4}
Let $k,n$ be integers with $1\le k\le n$, and let $a,b$ be positive integers satisfying
$2a\le b$. Set $\varepsilon=a/b$. The density on the sign cube obtained by transporting
the deterministic small-bias density through the almost $k$-wise transform has Fourier
coefficient of absolute value at most $\varepsilon$ on every nonempty set of at most $k$
coordinates.
\end{theorem}

\begin{theorem}[Output multiset cardinality bound in real scale]
\lean{FABL.almostKWiseOutputMultiset_card_le_realScale}
\proofsource{Analysis-of-Boolean-Functions-by-Ryan-ODonnell}{pdf-ffbe9c7b3306-p155-b007}
\label{FABL.almostKWiseOutputMultiset_card_le_realScale}
\leanok
\uses{FABL.AlmostKWiseInput, FABL.AlmostKWiseInput.epsilon, FABL.AlmostKWiseInput.epsilon_pos, FABL.AlmostKWiseInput.rowCount, FABL.AlmostKWiseInput.rowCount_le_sizeScale, FABL.AlmostKWiseInput.sizeScale, FABL.AlmostKWiseInput.smallBiasInput, FABL.almostKWiseOutputMultiset, FABL.almostKWiseOutputMultiset_card_eq_source, FABL.deterministicSmallBiasAlgorithm_spec, FABL.deterministicSmallBiasMultiset}
\difficulty{4}
Let an almost-wise input have independence order $k$ and output dimension $n$, with
$1\le k\le n$. Let its bias be the positive rational number whose numerator is $a$ and
denominator is $b$, where $2a\le b$, and write $\varepsilon=a/b$. The real-valued
cardinality of its output multiset is at most $16(S/\varepsilon)^2$, where $S$ is the
size scale.
\end{theorem}

\begin{theorem}[Output multiset size bound]
\lean{FABL.almostKWiseOutputMultiset_card_le_smallBiasScale}
\proofsource{Analysis-of-Boolean-Functions-by-Ryan-ODonnell}{
  pdf-ffbe9c7b3306-p155-b007,
  pdf-ffbe9c7b3306-p153-b002
}
\label{FABL.almostKWiseOutputMultiset_card_le_smallBiasScale}
\leanok
\uses{FABL.AlmostKWiseInput, FABL.AlmostKWiseInput.smallBiasInput, FABL.SmallBiasInput.fieldDegree, FABL.SmallBiasInput.fieldSize_le_two_scale, FABL.SmallBiasInput.scale, FABL.almostKWiseOutputMultiset, FABL.almostKWiseOutputMultiset_card_eq_source, FABL.deterministicSmallBiasMultiset_card}
\difficulty{3}
Let $k,n,a,b$ be positive integers with $k\le n$ and $2a\le b$. Form the almost $k$-wise
construction with output dimension $n$ and requested bias $a/b$, and let $R$ be its row
count and $S=\lceil Rb/a\rceil$. The cardinality of its output multiset is at most
$4S^2$.
\end{theorem}

\begin{theorem}[Output multiset cardinality bound]
\lean{FABL.almostKWiseOutputMultiset_card_le_naturalScale}
\proofsource{Analysis-of-Boolean-Functions-by-Ryan-ODonnell}{
  pdf-ffbe9c7b3306-p155-b007,
  pdf-ffbe9c7b3306-p170-b003
}
\label{FABL.almostKWiseOutputMultiset_card_le_naturalScale}
\leanok
\uses{FABL.AlmostKWiseInput, FABL.AlmostKWiseInput.outputScale, FABL.AlmostKWiseInput.smallBiasScale_le_outputScale, FABL.almostKWiseOutputMultiset, FABL.almostKWiseOutputMultiset_card_le_smallBiasScale}
\difficulty{1}
Let an almost-wise input have independence order $k$ and output dimension $n$, with
$1\le k\le n$. Let its bias be the positive rational number whose numerator is $a$ and
denominator is $b$, where $2a\le b$. The cardinality of its output multiset is at most
four times the square of its natural-valued output scale.
\end{theorem}

\begin{theorem}[Asymptotic output size in real scale]
\lean{FABL.almostKWiseOutputMultiset_card_isBigO_realScale}
\proofsource{Analysis-of-Boolean-Functions-by-Ryan-ODonnell}{
  pdf-ffbe9c7b3306-p155-b007,
  pdf-ffbe9c7b3306-p170-b003
}
\label{FABL.almostKWiseOutputMultiset_card_isBigO_realScale}
\leanok
\uses{FABL.AlmostKWiseInput, FABL.AlmostKWiseInput.epsilon, FABL.AlmostKWiseInput.sizeScale, FABL.almostKWiseOutputMultiset, FABL.almostKWiseOutputMultiset_card_le_realScale}
\difficulty{3}
Let an almost-wise input have independence order $k$ and output dimension $n$, with
$1\le k\le n$. Let its bias be the positive rational number whose numerator is $a$ and
denominator is $b$, where $2a\le b$, and write $\varepsilon=a/b$. As the real-valued
scale $S/\varepsilon$ tends to infinity, where $S$ is the size scale, the real-valued
cardinality of the output multiset is $O((S/\varepsilon)^2)$.
\end{theorem}

\begin{theorem}[Asymptotic output multiset size]
\lean{FABL.almostKWiseOutputMultiset_card_isBigO_naturalScale}
\proofsource{Analysis-of-Boolean-Functions-by-Ryan-ODonnell}{pdf-ffbe9c7b3306-p155-b007}
\label{FABL.almostKWiseOutputMultiset_card_isBigO_naturalScale}
\leanok
\uses{FABL.AlmostKWiseInput, FABL.AlmostKWiseInput.outputScale, FABL.almostKWiseOutputMultiset, FABL.almostKWiseOutputMultiset_card_le_naturalScale}
\difficulty{3}
Let an almost-wise input have independence order $k$ and output dimension $n$, with
$1\le k\le n$. Let its bias be the positive rational number whose numerator is $a$ and
denominator is $b$, where $2a\le b$. As the natural-valued output scale tends to
infinity, the real-valued cardinality of the output multiset is $O$ of the square of
that scale.
\end{theorem}

\begin{theorem}[Binary logarithm of a product]
\lean{FABL.binaryClog_mul_le}
\label{FABL.binaryClog_mul_le}
\leanok
\difficulty{3}
For all natural numbers $a$ and $b$, $\operatorname{clog}_2(ab) \le
\operatorname{clog}_2(a)+\operatorname{clog}_2(b)$, where $\operatorname{clog}_2$
denotes the ceiling of the base-two logarithm, with its natural-number convention at
zero.
\end{theorem}

\begin{definition}[Random-bit logarithmic scale]
\lean{FABL.AlmostKWiseInput.randomBitLogScale}
\statementsource{Analysis-of-Boolean-Functions-by-Ryan-ODonnell}{
  pdf-ffbe9c7b3306-p155-b006,
  pdf-ffbe9c7b3306-p170-b003
}
\label{FABL.AlmostKWiseInput.randomBitLogScale}
\leanok
\uses{FABL.AlmostKWiseInput, FABL.AlmostKWiseInput.reciprocalScale, FABL.AlmostKWiseInput.vandermondeInput, FABL.ExecutableVandermondeInput.fieldDegree}
Let $k,n$ be positive integers with $k\le n$, and let $a,b$ be positive integers with
$2a\le b$. The random-bit logarithmic scale is the sum of $\lceil\log_2 k\rceil$,
$\lceil\log_2(d+1)\rceil$, and $\lceil\log_2 s\rceil$, where $d=\max(1,\lceil\log_2
n\rceil)$ and $s$ is the reciprocal scale associated with the input.
\end{definition}

\begin{theorem}[Almost k wise random bits le log scale]
\lean{FABL.almostKWiseRandomBits_le_logScale}
\statementsource{Analysis-of-Boolean-Functions-by-Ryan-ODonnell}{pdf-ffbe9c7b3306-p155-b006}
\label{FABL.almostKWiseRandomBits_le_logScale}
\leanok
\uses{FABL.AlmostKWiseInput, FABL.AlmostKWiseInput.outputScale, FABL.AlmostKWiseInput.randomBitLogScale, FABL.AlmostKWiseInput.reciprocalScale, FABL.AlmostKWiseInput.sizeScale, FABL.AlmostKWiseInput.smallBiasInput, FABL.AlmostKWiseInput.smallBiasScale_le_outputScale, FABL.AlmostKWiseInput.vandermondeInput, FABL.ExecutableVandermondeInput.fieldDegree, FABL.SmallBiasInput.fieldDegree_eq_clog, FABL.SmallBiasInput.scale, FABL.almostKWiseRandomBits, FABL.binaryClog_mul_le}
\difficulty{4}
% dataset-pipeline: stub -- prose pending build_informal_statements.py
Statement pending regeneration.
\end{theorem}

\begin{theorem}[Almost k wise random bits is big o]
\lean{FABL.almostKWiseRandomBits_isBigO}
\statementsource{Analysis-of-Boolean-Functions-by-Ryan-ODonnell}{
  pdf-ffbe9c7b3306-p155-b006,
  pdf-ffbe9c7b3306-p170-b003
}
\label{FABL.almostKWiseRandomBits_isBigO}
\leanok
\uses{FABL.AlmostKWiseInput, FABL.AlmostKWiseInput.randomBitLogScale, FABL.almostKWiseRandomBits, FABL.almostKWiseRandomBits_le_logScale}
\difficulty{3}
% dataset-pipeline: stub -- prose pending build_informal_statements.py
Statement pending regeneration.
\end{theorem}

\begin{definition}[Executable matrix-vector multiplication work]
\lean{FABL.executableMatrixVecMulWork}
\label{FABL.executableMatrixVecMulWork}
\leanok
\uses{FABL.binaryConstructorTraversalWork}
For natural numbers $m$ and $n$, define the executable matrix-vector multiplication work
to be the binary-constructor traversal work with parameter $2m+1$ and the list of
indices from $0$ through $n-1$.
\end{definition}

\begin{theorem}[Matrix vector multiplication work]
\lean{FABL.executableMatrixVecMulWork_eq}
\label{FABL.executableMatrixVecMulWork_eq}
\leanok
\uses{FABL.binaryConstructorTraversalWork_eq, FABL.executableMatrixVecMulWork}
\difficulty{1}
For all natural numbers $m$ and $n$, the executable matrix-vector multiplication work is
$n(2m+1)$.
\end{theorem}

\begin{definition}[Work of the almost-wise construction]
\lean{FABL.almostKWiseConstructionWork}
\statementsource{Analysis-of-Boolean-Functions-by-Ryan-ODonnell}{
  pdf-ffbe9c7b3306-p155-b007,
  pdf-ffbe9c7b3306-p170-b003
}
\label{FABL.almostKWiseConstructionWork}
\leanok
\uses{FABL.AlmostKWiseInput, FABL.AlmostKWiseInput.rowCount, FABL.AlmostKWiseInput.smallBiasInput, FABL.binaryConstructorTraversalWork, FABL.deterministicSmallBiasAlgorithm, FABL.deterministicSmallBiasWork, FABL.executableMatrixVecMulWork, FABL.executableVandermondeConstructionWork}
For an input $I$ consisting of integers $k,n,a,b$ with $1\le k\le n$, $0<a$, and $2a\le
b$, the construction work is the sum of the executable Vandermonde construction work for
$k,n$, the deterministic small-bias work for the small-bias input associated with $I$,
and the binary-constructor traversal work evaluated at the executable matrix-vector
multiplication work for the row count and $n$, increased by one, and the number of
outputs of the deterministic small-bias algorithm for that input.
\end{definition}

\begin{theorem}[Work of the Almost-Wise Construction]
\lean{FABL.almostKWiseConstructionWork_eq}
\label{FABL.almostKWiseConstructionWork_eq}
\leanok
\uses{FABL.AlmostKWiseInput, FABL.AlmostKWiseInput.rowCount, FABL.AlmostKWiseInput.smallBiasInput, FABL.almostKWiseConstructionWork, FABL.almostKWiseRandomBits, FABL.almostKWiseSmallBiasOutputList_length, FABL.binaryConstructorTraversalWork_eq, FABL.deterministicSmallBiasWork, FABL.executableMatrixVecMulWork, FABL.executableVandermondeConstructionWork}
\difficulty{1}
For an almost-wise input with independence order $k$, output dimension $n$, bias
numerator $a$, and bias denominator $b$, the total construction work is the Vandermonde
construction work plus the deterministic small-bias construction work, plus $2^B$ times
the matrix-vector multiplication work per output, where $B$ is twice the field degree of
the associated small-bias input.
\end{theorem}

\begin{theorem}[Almost k wise vandermonde work le]
\lean{FABL.almostKWiseVandermondeWork_le}
\statementsource{Analysis-of-Boolean-Functions-by-Ryan-ODonnell}{
  pdf-ffbe9c7b3306-p154-b007,
  pdf-ffbe9c7b3306-p170-b003
}
\label{FABL.almostKWiseVandermondeWork_le}
\leanok
\uses{FABL.AlmostKWiseInput, FABL.AlmostKWiseInput.n_le_scale, FABL.AlmostKWiseInput.scale, FABL.AlmostKWiseInput.vandermondeInput, FABL.ExecutableVandermondeInput.fieldDegree, FABL.ExecutableVandermondeInput.fieldDegree_succ_le, FABL.ExecutableVandermondeInput.fieldSize_le_two_n, FABL.binaryAddWork, FABL.binaryArithmeticWork_le, FABL.binaryFieldPreprocessingWork, FABL.binaryFieldPreprocessingWork_le, FABL.binaryMulModWork, FABL.executableVandermondeConstructionWork, FABL.executableVandermondeConstructionWork_eq, FABL.executableVandermondePowerRowsWork, FABL.executableVandermondePowerRowsWork_le}
\difficulty{5}
% dataset-pipeline: stub -- prose pending build_informal_statements.py
Statement pending regeneration.
\end{theorem}

\begin{theorem}[Work bound for matrix–vector multiplication]
\lean{FABL.almostKWiseMatrixVecMulWork_le}
\label{FABL.almostKWiseMatrixVecMulWork_le}
\leanok
\uses{FABL.AlmostKWiseInput, FABL.AlmostKWiseInput.n_le_scale, FABL.AlmostKWiseInput.rowCount, FABL.AlmostKWiseInput.rowCount_le_four_scale_sq, FABL.AlmostKWiseInput.scale, FABL.executableMatrixVecMulWork, FABL.executableMatrixVecMulWork_eq}
\difficulty{4}
Let $k,n,a,b$ be natural numbers satisfying $1\le k\le n$, $a>0$, and $2a\le b$. Set
$s=\lceil nb/a\rceil$, and let $W$ be the work of the executable matrix–vector
multiplication with $n$ columns and $(k-1)\max(1,\lceil\log_2 n\rceil)+1$ rows, where
each traversal step costs twice the row count plus one. Then $W\le 9(s+1)^3$.
\end{theorem}

\begin{definition}[Polynomial budget]
\lean{FABL.AlmostKWiseInput.polynomialBudget}
\statementsource{Analysis-of-Boolean-Functions-by-Ryan-ODonnell}{
  pdf-ffbe9c7b3306-p155-b007,
  pdf-ffbe9c7b3306-p170-b003
}
\label{FABL.AlmostKWiseInput.polynomialBudget}
\leanok
\uses{FABL.AlmostKWiseInput, FABL.AlmostKWiseInput.scale}
For an input with scale $s$, define its polynomial budget to be $2^{38}(s+1)^{16}$.
\end{definition}

\begin{theorem}[Polynomial construction-work bound]
\lean{FABL.almostKWiseConstructionWork_le_polynomialBudget}
\label{FABL.almostKWiseConstructionWork_le_polynomialBudget}
\leanok
\uses{FABL.AlmostKWiseInput, FABL.AlmostKWiseInput.polynomialBudget, FABL.AlmostKWiseInput.rowCount, FABL.AlmostKWiseInput.scale, FABL.AlmostKWiseInput.smallBiasInput, FABL.AlmostKWiseInput.smallBiasScale_le_four_scale_sq, FABL.SmallBiasInput.polynomialBudget, FABL.SmallBiasInput.scale, FABL.almostKWiseConstructionWork, FABL.almostKWiseConstructionWork_eq, FABL.almostKWiseMatrixVecMulWork_le, FABL.almostKWiseOutputMultiset, FABL.almostKWiseOutputMultiset_card_eq_two_pow_randomBits, FABL.almostKWiseOutputMultiset_card_le_smallBiasScale, FABL.almostKWiseRandomBits, FABL.almostKWiseVandermondeWork_le, FABL.deterministicSmallBiasWork, FABL.deterministicSmallBiasWork_le_polynomialBudget, FABL.executableMatrixVecMulWork, FABL.executableVandermondeConstructionWork}
\difficulty{5}
For every input to the almost-wise independent construction, its total construction work
is at most $2^{38}(s+1)^{16}$, where $s=\lceil nb/a\rceil$ is the input scale for bias
$a/b$.
\end{theorem}

\begin{theorem}[Asymptotic construction-work bound]
\lean{FABL.almostKWiseConstructionWork_isBigO}
\proofsource{Analysis-of-Boolean-Functions-by-Ryan-ODonnell}{
  pdf-ffbe9c7b3306-p155-b007,
  pdf-ffbe9c7b3306-p170-b003
}
\label{FABL.almostKWiseConstructionWork_isBigO}
\leanok
\uses{FABL.AlmostKWiseInput, FABL.AlmostKWiseInput.polynomialBudget, FABL.AlmostKWiseInput.scale, FABL.almostKWiseConstructionWork, FABL.almostKWiseConstructionWork_le_polynomialBudget}
\difficulty{3}
For inputs to the almost-wise independent construction, let $s=\lceil nb/a\rceil$, where
$a/b$ is the input bias. As $s$ tends to infinity over the admissible inputs, the total
construction work is $O((s+1)^{16})$.
\end{theorem}

\begin{theorem}[Guarantees of the almost-wise independent construction]
\lean{FABL.almostKWiseAlgorithm_spec}
\proofsource{Analysis-of-Boolean-Functions-by-Ryan-ODonnell}{
  pdf-ffbe9c7b3306-p154-b006,
  pdf-ffbe9c7b3306-p154-b007,
  pdf-ffbe9c7b3306-p155-b004,
  pdf-ffbe9c7b3306-p155-b007
}
\label{FABL.almostKWiseAlgorithm_spec}
\leanok
\uses{FABL.AlmostKWiseInput, FABL.AlmostKWiseInput.epsilon, FABL.AlmostKWiseInput.epsilon_le_half, FABL.AlmostKWiseInput.epsilon_pos, FABL.AlmostKWiseInput.outputScale, FABL.AlmostKWiseInput.polynomialBudget, FABL.AlmostKWiseInput.randomBitLogScale, FABL.AlmostKWiseInput.rowCount, FABL.AlmostKWiseInput.scale, FABL.AlmostKWiseInput.sizeScale, FABL.AlmostKWiseInput.smallBiasInput, FABL.IsLowDegreeFourierRegular, FABL.almostKWiseConstructionWork, FABL.almostKWiseConstructionWork_eq, FABL.almostKWiseConstructionWork_isBigO, FABL.almostKWiseConstructionWork_le_polynomialBudget, FABL.almostKWiseDensity, FABL.almostKWiseDensity_isApproximatelyKWiseIndependent, FABL.almostKWiseOutputList, FABL.almostKWiseOutputList_length, FABL.almostKWiseOutputMultiset, FABL.almostKWiseOutputMultiset_card_eq_two_pow_randomBits, FABL.almostKWiseOutputMultiset_card_isBigO_naturalScale, FABL.almostKWiseOutputMultiset_card_isBigO_realScale, FABL.almostKWiseOutputMultiset_card_le_naturalScale, FABL.almostKWiseOutputMultiset_card_le_realScale, FABL.almostKWiseRandomBits, FABL.almostKWiseRandomBits_isBigO, FABL.almostKWiseRandomBits_le_logScale, FABL.binaryFunctionOnSignCube, FABL.deterministicSmallBiasWork, FABL.executableMatrixVecMulWork, FABL.executableVandermondeConstructionWork}
\difficulty{1}
Let $k,n,a,b$ be natural numbers with $1\le k\le n$, $a>0$, and $2a\le b$. Set
$d=\max(1,\lceil\log_2 n\rceil)$, $r=\lceil b/a\rceil$, $s=k(d+1)$, $q=sr$, and
$t=\lceil nb/a\rceil$, where the ceilings of rational quantities are rounded up to
natural numbers. For the construction with these parameters, let $L$ be its output list,
let $M$ be the cardinality of its output multiset, let $W$ be its construction work, and
let $B=2d$ be its random-bit count. Then $0<a/b\le 1/2$, $|L|=M=2^B$,
$M\le16(s/(a/b))^2$, and $M\le4q^2$. As $s/(a/b)$ tends to infinity along such inputs,
$M=O((s/(a/b))^2)$; as $q$ tends to infinity, $M=O(q^2)$. The construction's probability
density, transferred to the sign cube, has every nonempty Fourier coefficient on at most
$k$ coordinates bounded in absolute value by $a/b$. Also, $B\le2(\lceil\log_2
k\rceil+\lceil\log_2(d+1)\rceil+\lceil\log_2 r\rceil)$, and $B$ is $O(\lceil\log_2
k\rceil+\lceil\log_2(d+1)\rceil+\lceil\log_2 r\rceil)$ as that quantity tends to
infinity. The work is the sum of the Vandermonde construction work and the small-bias
construction work, plus $2^B$ times the matrix-vector multiplication work per output; it
is at most $2^{38}(t+1)^{16}$ and is $O((t+1)^{16})$ as $t$ tends to infinity.
\end{theorem}
